% ── APPENDIX A: FULL RCT PROOF ───────────────────────────────────────────────

\section{Full Proof of the Relational Coherence Theorem}
\label{app:rct-proof}

\subsection{Mathematical Preliminaries}

Let $\mathcal{H}_S$ and $\mathcal{H}_E$ be finite-dimensional Hilbert spaces
with $\dim\mathcal{H}_S = d_S \geq 2$ and $\dim\mathcal{H}_E = d_E \geq d_S$.
The composite system occupies $\mathcal{H} = \mathcal{H}_S \otimes \mathcal{H}_E$.
Observables are self-adjoint elements of the C*-algebra
$\mathcal{B}(\mathcal{H})$ of bounded linear operators on $\mathcal{H}$.
States are density operators: positive trace-class operators
$\rho \in \mathcal{B}(\mathcal{H})$ with $\mathrm{Tr}[\rho] = 1$.

Fix orthonormal bases $\{\ket{s_i}\}_{i=1}^{d_S}$ and
$\{\ket{e_k}\}_{k=1}^{d_E}$ for $\mathcal{H}_S$ and $\mathcal{H}_E$
respectively.
Any pure state of $SE$ admits the Schmidt decomposition
\begin{equation}
  \ket{\psi_{SE}} = \sum_{j=1}^{r} \sqrt{\lambda_j}\,\ket{u_j}\otimes\ket{v_j},
  \label{eq:schmidt}
\end{equation}
where $r \leq \min(d_S, d_E)$ is the Schmidt rank, $\lambda_j > 0$,
$\sum_j \lambda_j = 1$, and $\{\ket{u_j}\}$, $\{\ket{v_j}\}$ are
orthonormal sets in $\mathcal{H}_S$ and $\mathcal{H}_E$ respectively.
The state is a product state iff $r = 1$; it is maximally entangled iff
$r = \min(d_S, d_E)$ and all $\lambda_j = 1/r$.

\subsection{Proof of Proposition~\ref{prop:indefiniteness} (Property Indefiniteness)}

\textbf{Claim.} There is no function $\lambda: \mathcal{H}_S \to \mathrm{spec}(\hat{A})$
assigning definite eigenvalues to all observables $\hat{A}$ on $\mathcal{H}_S$
prior to relational interaction, consistently with all quantum predictions.

\textbf{Proof.}
We treat the two cases $d_S = 2$ and $d_S \geq 3$ separately.

\textit{Case $d_S = 2$ (qubit).}
Suppose, for contradiction, that definite pre-existing values
$\lambda(A) \in \{+1,-1\}$ exist for all qubit observables.
For any product state $\ket{\psi_S}\otimes\ket{\psi_E}$ of two qubits
and measurement directions $\mathbf{a}, \mathbf{a}', \mathbf{b}, \mathbf{b}'$,
the CHSH combination satisfies
\begin{equation}
  \bigl|\langle\hat{A}\hat{B}\rangle - \langle\hat{A}\hat{B}'\rangle
    + \langle\hat{A}'\hat{B}\rangle + \langle\hat{A}'\hat{B}'\rangle\bigr|
  \leq 2.
\end{equation}
For the maximally entangled state $\ket{\Phi^+}$ with directions chosen at
$45^\circ$ intervals, quantum mechanics predicts a violation of $2\sqrt{2}$,
confirmed experimentally~\cite{Aspect1982,Zeilinger2023Nobel}.
This contradiction establishes that no pre-existing value assignment exists
for $d_S = 2$.

\textit{Case $d_S \geq 3$.}
The Kochen--Specker theorem~\cite{KochenSpecker1967} proves that for
$d_S \geq 3$, there is no non-contextual hidden variable assignment
consistent with the quantum predictions for all projection-valued
measures on $\mathcal{H}_S$.
(A non-contextual assignment would give the same value to a projector
$\hat{P}$ regardless of which complete set of commuting observables it
is measured alongside~---~precisely the assignment required by any
hidden variable theory seeking to recover quantum predictions.)
No such assignment exists.

In both cases, properties are not possessed independently of measurement
context, i.e., independently of relational interaction. $\square$

\subsection{Proof of Proposition~\ref{prop:determination} (Relational Determination)}

\textbf{Claim.} A definite outcome for observable $\hat{A}$ on $S$ occurs
if and only if $S$ undergoes a relational interaction with some $E_i$ that
increases $\bar\alpha(S, E_i)$ from zero.

\textbf{Proof.}

($\Rightarrow$, \emph{necessity})
Suppose $S$ is relationally isolated: $\bar\alpha(S, E_i) = 0$ for all $i$.
Then $\|\alpha(S,E_i)\|_F = 0$, so all off-diagonal elements of
$\rhoSE$ vanish and $\rhoSE$ is diagonal.
For a pure state $\ket{\psi_S}$, a diagonal reduced density matrix requires
$\ket{\psi_{SE}} = \ket{\psi_S}\otimes\ket{e_0}$ (a product state, Schmidt
rank 1).
The state $\ket{\psi_S}$ is a superposition of eigenstates of $\hat{A}$
unless it happens to be an eigenstate~---~but by Proposition~\ref{prop:indefiniteness},
no eigenstate is preferred prior to interaction.
Hence no definite outcome obtains.

($\Leftarrow$, \emph{sufficiency})
Suppose $S$ interacts with $E_i$ via Hamiltonian $H_{SE_i}$, generating
the joint unitary $U(t) = e^{-iH_{SE_i}t}$.
By the Schmidt decomposition, the evolved state
\begin{equation}
  \ket{\psi_{SE_i}(t)} = U(t)\ket{\psi_S}\otimes\ket{e_0}
\end{equation}
acquires Schmidt rank $r > 1$ for generic $H_{SE_i}$ and $t > 0$,
so $\bar\alpha(S, E_i) > 0$.
Zurek's einselection~\cite{Zurek2003} selects the pointer basis
$\{\ket{s_i^*}\}$ determined by $H_{SE_i}$, in which the reduced state
$\rhoSE$ becomes effectively diagonal.
Relative to $E_i$, $S$ has a definite outcome in the pointer basis. $\square$

\subsection{Proof of Proposition~\ref{prop:correspondence} (Coherence--Affinity Correspondence)}

\textbf{Claim.} $\bar\alpha = 0$ iff $\rhoSE$ is separable;
$\bar\alpha = 1$ iff $\ket{\psi_{SE}}$ is maximally entangled;
for pure states $\bar\alpha = \mathcal{C}(\rho_{SE})$.

\textbf{Proof.}

($\bar\alpha = 0 \Leftrightarrow$ separable)
$\bar\alpha = 0$ iff $\|\alpha\|_F = 0$ iff all off-diagonal elements of
$\rhoSE$ vanish iff $\rhoSE$ is diagonal in the Schmidt basis iff
$\ket{\psi_{SE}}$ has Schmidt rank 1 (product state).
For a product state, the concurrence $\mathcal{C} = 0$ and entanglement entropy $S = 0$.

($\bar\alpha = 1 \Leftrightarrow$ maximally entangled)
$\bar\alpha = 1$ iff $\|\alpha\|_F = \|\alpha_{\max}\|_F$.
For a two-qubit pure state in the Schmidt basis
$\ket{\psi_{SE}} = \sqrt{\lambda}\ket{u_1 v_1} + \sqrt{1-\lambda}\ket{u_2 v_2}$,
the off-diagonal element magnitude is $|\rho_{12}| = \sqrt{\lambda(1-\lambda)}$,
maximized at $\lambda = 1/2$ (i.e., $\mathcal{C} = 2\sqrt{\lambda(1-\lambda)} = 1$).
Thus $\bar\alpha = 1$ iff $\lambda = 1/2$ iff $\ket{\psi_{SE}}$ is a Bell state.

(Pure state identity $\bar\alpha = \mathcal{C}$)
For a two-qubit pure state with Schmidt coefficients
$\sqrt{\lambda}, \sqrt{1-\lambda}$:
$\mathcal{C} = 2\sqrt{\lambda(1-\lambda)}$ and
$\bar\alpha = |\rho_{12}|/|\rho_{12}^{\max}| = 2\sqrt{\lambda(1-\lambda)}$.
Hence $\bar\alpha = \mathcal{C}$.

(Mixed state inequality $\bar\alpha \leq \mathcal{C}$)
Follows from the concavity of the Frobenius norm of the coherences under
convex mixture: for $\rho = \sum_k p_k \rho_k$,
$\|\alpha(\rho)\|_F \leq \sum_k p_k\|\alpha(\rho_k)\|_F$,
while $\mathcal{C}(\rho) = \max\sum_k p_k \mathcal{C}(\rho_k)$ by convex
roof construction~\cite{Wootters1998}. $\square$

\subsection{Proof of Proposition~\ref{prop:conservation} (Conservation of Relational Potential)}

\textbf{Claim.} For a closed system with unitary evolution $U(t)$,
$\mathrm{Tr}[\rho(t)^2] = \mathrm{Tr}[\rho(0)^2]$ for all $t$.

\textbf{Proof.}
\begin{align}
  \mathrm{Tr}\bigl[\rho(t)^2\bigr]
    &= \mathrm{Tr}\bigl[U(t)\rho(0)U^\dagger(t)\cdot U(t)\rho(0)U^\dagger(t)\bigr] \\
    &= \mathrm{Tr}\bigl[U(t)\rho(0)^2 U^\dagger(t)\bigr]
      \qquad (\text{since } U^\dagger U = \mathbf{I}) \\
    &= \mathrm{Tr}\bigl[\rho(0)^2\bigr]
      \qquad (\text{cyclicity of trace}).
\end{align}
The apparent purity loss in the reduced state
$\rhoSE = \mathrm{Tr}_E[\rho_{SE}]$ reflects transfer of coherence
into $SE$ correlations~---~the joint state remains pure, and
$\mathrm{Tr}[\rho_{SE}^2] = 1$ throughout.
The ``lost'' coherence is not destroyed; it is encoded in the relational
network $\Gamma_S$ as $SE$ entanglement. $\square$

\subsection{Weyl Chamber Coordinates for $\URA(\alpha)$}

The local invariants of a two-qubit gate $U \in \mathrm{SU}(4)$ are
computed from the matrix $M = U^T_B U_B$ in the Bell basis, where
$U_B = \frac{1}{\sqrt{2}}\begin{pmatrix}1&0&0&i\\0&i&1&0\\0&i&{-1}&0\\1&0&0&{-i}\end{pmatrix}$~\cite{Zhang2003WeylChamber}.

For $\URA(\alpha)$ defined by Eq.~\eqref{eq:URA}, direct computation gives
$M = \mathrm{diag}(e^{-i\alpha}, e^{i\alpha}, e^{i\alpha}, e^{-i\alpha})$
in the Bell basis, yielding Weyl coordinates:
\begin{equation}
  (c_1, c_2, c_3) = (\alpha, \alpha, \alpha), \quad \alpha \in [0, \pi/2].
  \label{eq:weyl-full}
\end{equation}
This trajectory lies entirely within the \textsc{iswap} family of gates,
confirming that $\URA$ is an isotropic exchange gate continuously
interpolating from the identity ($\alpha = 0$) to the \textsc{swap}
gate ($\alpha = \pi/2$), with maximum entanglement capability at
$\alpha = \pi/4$ (the \textsc{iswap} point).

\subsection{Simulation Parameters}
\label{app:sim-params}

All simulations use the following configuration unless otherwise noted.

\textbf{Qiskit (IBM):} \texttt{qiskit-terra} v0.45, \texttt{qiskit-aer} v0.13.
Two-qubit system; $10^4$ shots per circuit; state vector simulation for
process fidelity; density matrix simulation for noise models.
Gate decomposition: $\URA(\alpha)$ compiled via \texttt{transpile} with
optimization level~3 to the $\{R_Z, \sqrt{X}, \textsc{cx}\}$ native gate set.

\textbf{Cirq (Google):} \texttt{cirq-core} v1.3.
Same qubit configuration; \texttt{DensityMatrixSimulator} with noise models
applied as \texttt{ConstantQubitNoiseModel}.

\textbf{Noise models:}
\begin{enumerate}
  \item \emph{Depolarizing}: single-qubit error rate $p \in \{0.001, 0.005, 0.01, 0.02, 0.05\}$.
  \item \emph{Dephasing (phase damping)}: dephasing rate $\gamma \in \{0.01, 0.05, 0.1, 0.2\}$.
  \item \emph{Amplitude damping}: $T_1$ relaxation probability $p_{\rm AD} \in \{0.01, 0.05, 0.1\}$.
\end{enumerate}

\textbf{Process fidelity} is computed via quantum process tomography using
the Choi--Jamio\l{}kowski isomorphism, comparing the noisy process matrix
$\chi_{\rm noisy}$ to the ideal $\chi_{\rm ideal}$:
$F_{\rm proc} = \mathrm{Tr}[\chi_{\rm ideal}\,\chi_{\rm noisy}]$.

\textbf{Baseline gates:} \textsc{cnot} = $\mathrm{CX}$ gate at $(c_1,c_2,c_3) = (\pi/2,0,0)$;
\textsc{iswap} at $(\pi/4,\pi/4,0)$; \textsc{cz} at $(\pi/2,0,0)$ (locally
equivalent to \textsc{cnot}).
All baselines implemented via native Qiskit/Cirq gate objects.
